% J2.6: del dato a la decisión; cada flecha es una elección que hay que justificar.
\documentclass[tikz,border=4pt]{standalone}
\usepackage{resonancia}
\begin{document}
\begin{tikzpicture}
  \node[cajaAzul, minimum width=2.5cm] (q) at (0,0) {\textbf{pregunta}\\[-1pt]{\footnotesize ¿ha cambiado $f$?}};
  \node[cajaVerde, minimum width=2.5cm] (h) at (3.4,0) {\textbf{hipótesis}\\[-1pt]{\footnotesize $H_0$: no ha cambiado}};
  \node[cajaNaranja, minimum width=2.5cm] (e) at (6.8,0) {\textbf{contraste}\\[-1pt]{\footnotesize $t$, $p$, IC 95 \%}};
  \node[caja, draw=rVioleta, fill=rVioleta!8, minimum width=2.5cm] (r) at (10.2,0) {\textbf{relevancia}\\[-1pt]{\footnotesize tamaño de efecto}};
  \draw[flecha] (q) -- (h); \draw[flecha] (h) -- (e); \draw[flecha] (e) -- (r);
  \node[etiqueta, color=rGris, anchor=north] at (5.1,-0.9) {supuestos: independencia, normalidad aproximada, comparaciones declaradas de antemano};
\end{tikzpicture}
\end{document}
